\documentclass[10pt]{article}

\usepackage[utf8]{inputenc}

\usepackage[T1]{fontenc}

\usepackage{amsmath}

\usepackage{amsfonts}

\usepackage{amssymb}

\usepackage[version=4]{mhchem}

\usepackage{stmaryrd}

\usepackage{CJKutf8}

\usepackage{bbold}



\begin{document}

КОШИ ИНТЕГРАЛЬНАЯ ФОРМУЛА - см. Аналитическая функция.



КОШИ ПОВЕРХНОСТЬ - см. Пространство-время теории относительности.



КОШИ TEOPEMA - теорема об обрашении в нуль интеграла от аналитической функции, взятого вдоль замкнутого контура. Точнее, пусть функция $f(t)$ аналитична в области $D$, а $\gamma$ - кусочно гладкий контур, лежащий в $D$ и не содержащий внутри себя особенностей функции $f(z)$. Тогда, согласно К. т., контурный интеграл



\[

\int_{\gamma} f(z) d z

\]



равен нулю. Доказана О. Коши (A. Cauchy, 1825). Геометрически К. т. означает, что векторное поле, компонентами к-рого являются, соответственно, вещественная и мнимая части аналитич. функции, потенциально и соленоидально, то есть его дивергенция и ротор равны нулю. Справедливо и обратное утверждение (т е о р е ма М о р е ры): если функция $f(z)$ непрерывна в односвязной области $D$ и такова, что



\[

\int_{\gamma} f(z) d z=0

\]



для любого кусочно гладкого замкнутого контура $\gamma$, лежащего в $D$, то $f(z)$ аналитична в $D$. К. т. играет важную роль в теории аналитич. функщий. На ней основано представление аналитич. функции в виде Коши интеграла, она используется в теории вычетов и т. Д.



См. также Дифференциальное уравнение; аналитическая теория.



КОШИ-АДАМАРА ТЕОРЕМА - см. СтепенНой ряд. КОШИ-АДАМАРА ФОРМУЛА - сМ. Степенной ряд. КОШИ-ПУАНКАРЕ ИНТЕГРАЛЬНАЯ TЕОРЕМА - см. Аналитическая функция.



КОШИ-РИМАНА СИСТЕМА - см. Аналитическая функция.



КОШи-РИмАнА условия, Д' Алам бе ра-\begin{CJK}{UTF8}{mj}字\end{CJK} ли е ра условия,- условия на действительную $u=u(x, y)$ и мнимую $v=v(x, y)$ части функции комплексного переменного



\[

w=f(z)=u+i v, z=x+i y,

\]



обеспечивающие моногенность и аналитичность $f(z)$ как функции комплексного переменного.



Для того чтобы функция $w=f(z)$, определенная в нек-рой области $D$ комплексной плоскости $z$, была м о н о ге н н о й в точке $z_{0}=x_{0}+i y_{0}$, то есть имела производную в точке $z_{0}$ как функция комплексного переменного $z$, необходимо и достаточно, чтобы ее действительная и мнимая части $\boldsymbol{u}$ и $\boldsymbol{v}$ были дифференцируемы в точке $\left(x_{0}, y_{0}\right)$ как функции действительных переменных $x$ и $y$ и чтобы, кроме того, в этой точке выполнялись условия Коши-Римана:



\[

\frac{\partial u}{\partial x}=\frac{\partial v}{\partial y}, \quad \frac{\partial u}{\partial y}=-\frac{\partial v}{\partial x} .

\]



Если К.-Р. у. выполнены, то производная $f^{\prime}(z)$ представима в любой из следующих форм:



\[

f^{\prime}(z)=\frac{\partial u}{\partial x}+i \frac{\partial v}{\partial x}=\frac{\partial v}{\partial y}-i \frac{\partial u}{\partial y}=\frac{\partial u}{\partial x}-i \frac{\partial u}{\partial y}=\frac{\partial v}{\partial y}+i \frac{\partial v}{\partial x} .

\]



Для того чтобы однозначная в области $D$ функция $f(z)$ была аналитической в $D$, необходимо и достаточно, чтобы ее действительная и мнимая части были дифференцируемыми функциями, удовлетворяющими К.-Р.у. всюду в $D$. Функции $u$ и $v$ класса $C^{2}(D)$, удовлетворяющие К.-Р. у. (1), являются, каждая по отдельности, гармоническими функциями от $x$ и $\boldsymbol{y}$; условия (1) суть усло в ия с опря жен ност и этих двух гармонич. функций: зная одну из них, вторую можно найти интегрированием. Условия (1) справедливы для любых двух ортогональных направлений $s$ и $n$, ориентированных взаимно так же, как оси $О \boldsymbol{x}$ и $О y$ в виде:



\[

\frac{\partial u}{\partial s}=\frac{\partial v}{\partial n}, \quad \frac{\partial u}{\partial n}=-\frac{\partial v}{\partial s}

\]



Напр., в полярных координатах $(r, \varphi)$ при $r \neq 0$ :



\[

\frac{\partial u}{\partial r}=\frac{1}{r} \frac{\partial v}{\partial \varphi}, \quad \frac{1}{r} \frac{\partial u}{\partial \varphi}=-\frac{\partial v}{\partial r} .

\]



Введя операторы комплексного дифферен цирования



\[

\frac{\partial}{\partial z}=\frac{1}{2}\left(\frac{\partial}{\partial x}-i \frac{\partial}{\partial y}\right), \frac{\partial}{\partial \bar{z}}=\frac{1}{2}\left(\frac{\partial}{\partial x}+i \frac{\partial}{\partial y}\right),

\]



К.-Р. у. (1) можно записать в виде



\[

\frac{\partial f}{\partial \bar{z}}=0

\]



Таким образом, дифференцируемая функция $f(z, \bar{z})$ переменных $z$ и $\bar{z}$ является аналитич. функцией от $z$ тогда и только тогда, когда $\partial f / \partial \bar{z}=0$.



Для аналитич. функций многих комплексных переменных



\[

z=\left(z_{1}, z_{2}, \ldots, z_{n}\right), z_{k}=x_{k}+i y_{k}, k=1,2, \ldots, n,

\]



К.-Р.у. записываются в виде переопределенной при $n>1$ системы уравнений с частными производными для функций



\[

\begin{aligned}

& u=\operatorname{Re} f(z)=u\left(x_{1}, x_{2}, \ldots, x_{n} ; y_{1}, y_{2}, \ldots, y_{n}\right), \\

& v=\operatorname{Im} f(z)=v\left(x_{1}, x_{2}, \ldots, x_{n} ; y_{1}, y_{2}, \ldots, y_{n}\right): \\

& \frac{\partial u}{\partial x_{k}}=\frac{\partial v}{\partial y_{k}}, \quad \frac{\partial u}{\partial y_{k}}=-\frac{\partial v}{\partial x_{k}}, \quad k=1,2, \ldots, n,

\end{aligned}

\]



или, при помощи операторов комплексного дифференцирования, в виде:



\[

\frac{\partial f}{\partial \bar{z}_{k}}=0, \quad k=1,2, \ldots, n .

\]



Функщии $u$ и $v$ класса $\mathbb{C}^{2}$, удовлетворяющие условиям (2), являются при $n \geqslant 1$, каждая по отдельности, плюригармоническими функциями от переменных $x_{k}$ и $y_{k}$. Плюригармонич. функции при $n>1$ составляют собственный подкласс класса гармонич. функций. Условия (2) суть усло в и я соп р я же н н ос т и двух плюригармонич. функций $u$ и $v$ : зная одну из них, другую можно найти интегрированием.



Лит.: [1] М а р к у ш е в и ч А. И., Теория аналитических функций, 2 изд., т. 1, М., 1967, гл. 1; [2] Ш а б а т Б. В., Введение в комплексный анализ, 2 изд., М., 1976, ч. 1, гл. 1, ч. 2, гл. 1.Е. Д. Соломенцев. КОШИ-ШВАРЦА НЕРАВЕНСТВО - сМ. Швариа неравенство.



КОЭФФИЦИЕНТНАЯ ФУНКЦИЯ - см. Произвоящий функционал.



KРАЕВАЯ ЗАДАЧА - задача, в которой из некоторого класса функций, определенных в данной области, требуется найти ту, которая удовлетворяет на границе (крае) этой области заданным условиям. Функции, описывающие конкретные явления природы (физич., химич. и др.), как правило, представляют собой решения уравнений математич. физики, выведенных из общих законов, к-рым подчиняются эти явления. Когда рассматриваемые уравнения допускают целые семейства решений, дополнительно задают так наз. краевые (граничные) или начальные услов и я, позволяюшие однозначно выделить интересующее нас решение. В то время как краевые условия задаются исключительно на граничных точках области, где ищется решение, начальные условия могут оказаться заданными на определенном множестве точек внутри области. Напр., уравнение



\[

\frac{\partial^{2} u}{\partial x_{1}^{2}}-\frac{\partial^{2} u}{\partial x_{2}^{2}}=0

\]



имеет бесконечное множество решений



\[

u\left(x_{1}, x_{2}\right)=f\left(x_{1}+x_{2}\right)+f_{1}\left(x_{1}-x_{2}\right),

\]





\end{document}